\subsection{Signed shuffle freeness and relative bar base change}

\begin{theorem}[C24: \status{Proved}; the relative bar property]
\label{sa08:C24}
The Koszul-signed shuffle algebra of a graded free module is a free
graded PD algebra. For an ordered homogeneous basis, its odd exterior
generators are the odd Lyndon words, and its even divided-power
generators are the even Lyndon words together with the concatenation
squares of odd Lyndon words. This description is compatible with
inclusions of ordered alphabets.

Let $M\hookrightarrow N$ be a split inclusion of free modules with
free quotient, and use the positive bar closures $U_r$ of
Statement~\ref{sa08:D029}. Then $U_r(N)$ is a semifree DG module
over $U_r(M)$, and $U_r(M)\to U_r(N)$ is a quasi-isomorphism.
For every strict PD differential algebra map $U_r(M)\to A$, the
ordinary PD pushout induces a quasi-isomorphism
\[
 A\longrightarrow A\otimes_{U_r(M)}U_r(N).
 \tag{C-24a}
\]
This conclusion is universal under coefficient extension and
preserves all additional weights. For a degreewise split inclusion
of free simplicial modules, it holds in every simplex.
\end{theorem}
\begin{proof}
We first derive the signed freeness assertion integrally. The primary
unsigned input is the precise triangularity theorem of
Melan\c{c}on--Reutenauer, \S4, printed pp.~589--590
\cite{C_melancon_reutenauer1989}: if the decreasing Lyndon
factorization of a word is
\[
 w=l_1^{n_1}\cdots l_t^{n_t},\qquad l_1>\cdots>l_t,
\]
the unsigned shuffle of these factors has leading word $w$ with
coefficient $\prod_i n_i!$, and every other word is smaller.
The signed result used here is proved next, rather than included
silently in that citation.

The signed shuffle algebra over $\Z$ has integral divided powers
on even positive elements. For an even basis word, the permutations
of its $a$ labelled shuffle blocks act freely and preserve the
sign, so its $a$-th shuffle power is divisible by $a!$. The sum
formula extends this to every even element. The PD identities
hold over $\Q$ and hence in the torsionfree integral word module.
Odd squares vanish by the sign-reversing block interchange.
All these integral operations extend to any coefficient ring.

For an even Lyndon factor $l$ of multiplicity $n$, use
$\gamma_n(l)$. For an odd factor of multiplicity $2m+\epsilon$,
where $\epsilon\in\{0,1\}$, use
\[
 \gamma_m(ll)\mathbin{\star}l^\epsilon,
\]
where $ll$ is concatenation and the last exponent is only zero or
one. Multiply these expressions in decreasing Lyndon order.
No term greater than $w$ can occur, by unsigned triangularity.
The shuffles producing $w$ are exactly permutations of whole
identical Lyndon blocks: the obvious such permutations already
account for the entire unsigned leading coefficient.

For the odd factor, the two copies within each concatenation
square must retain their order. Thus its signed leading coefficient
is the sign sum for shuffling $m$ ordered pairs of odd objects,
with at most one additional odd object. Write $B(a,b)$ for the
sign sum of $(a,b)$-shuffles of odd objects. The recursion
\[
 B(a,b)=B(a-1,b)+(-1)^aB(a,b-1),\qquad
 B(a,0)=B(0,b)=1,
\]
gives
\begin{align*}
 B(2a,2b)&=\binom{a+b}{a},\\
 B(2a+1,2b)=B(2a,2b+1)&=\binom{a+b}{a},\\
 B(2a+1,2b+1)&=0.
\end{align*}
The required count is consequently
$\prod_{j=1}^m B(2j-2,2)=m!$; the optional single block
contributes $B(2m,1)=1$. Dividing by $m!$ gives leading
coefficient one. Even Lyndon factors have leading coefficient
$n!$, canceled by their divided-power denominator.

The stated PD/exterior monomials therefore form a unitriangular
integral basis of the word module. Inverting the triangular
matrix is a finite operation on each fixed multiset of letters,
so the argument also applies to infinite alphabets. These are
exactly the monomials in the claimed free graded PD algebra.
An inclusion of ordered alphabets preserves all of its Lyndon
generators and their squares, proving compatibility.

Now choose a basis of $M$ and extend it to $N$. The divided-power
monomial alphabet of $\Gamma^+(M)$ is a subset of that for
$\Gamma^+(N)$. Applying the signed basis to the bar factors shows
that $U_r(N)$ is free as a graded PD $U_r(M)$-algebra. The extra
generators consist of the new coefficient generators and the new
signed-Lyndon bar generators. In particular it is free as a
graded module over $U_r(M)$.

Filter this free module by the total bar length of its new-generator
basis monomials, without counting factors in the old algebra
$U_r(M)$. The bar differential removes one letter, either by
moving the first letter into the coefficient or by multiplying
two adjacent letters. Expansion in the signed Lyndon basis
preserves total bar length. Hence the differential of a new
basis monomial is a sum of old coefficients times new basis
monomials of strictly smaller bar length: any old factor consumes
nonnegative length. The new coefficient generators have
differential zero. The filtration quotients are thus direct sums
of shifts of $U_r(M)$ with its own differential. This is a
semifree DG module filtration and proves K-flatness over every
coefficient ring.

Both bar closures augment quasi-isomorphically to $k$ by
Statement~\ref{sa08:D029}. Therefore $U_r(M)\to U_r(N)$ is a
quasi-isomorphism and its cone is an acyclic K-flat
$U_r(M)$-module. An acyclic K-flat module tensored with any
DG module is acyclic: choose a K-flat resolution of the latter,
then use the two K-flat properties to compare its tensor with
the original module and with its resolution. Tensoring the cone
therefore proves (C-24a).

The graded freeness just established identifies that tensor with
the ordinary PD pushout: after the map to $A$, the same extra
generators freely generate a graded PD $A$-algebra. Its
differential is induced by the bar differential. Every step
preserves weights and survives coefficient extension. For
simplicial modules, compatible basis choices are unnecessary:
the bar map is functorial, while an extending basis verifies
K-flatness separately in each simplex.
\end{proof}

\paragraph{Scope of the relative theorem.}
Statement~\ref{sa08:C24} proves the vertical base-change property
for a bar cell once its boundary map $U_r(M)\to A$ has been
given. A prescribed map on the original odd disk does not
automatically extend to this boundary map. Nor does the theorem
assert eventual horizontal bounds for arbitrary sequences of
relative bar attachments. Those are separate parts of the
remaining integral comparison problem.

\paragraph{Exact combinatorial check.}
The program \texttt{code/C\_signed\_shuffle\_checks.py} verifies
all $4956$ unitriangular columns for every parity assignment on
two letters through length eight and on three letters through
length five. It also checks integral divided-power coefficients
and vanishing odd squares. All checks pass; the complete receipt
is \texttt{logs/C\_signed\_shuffle\_checks.json}, with runtime
$0.083071112$ seconds. This corroborates the finite combinatorics;
the proof above covers arbitrary alphabets and lengths.
